% !TeX root = ../../Book2.tex
% 纲要标题：### 11.2 乘积、传递与自然性
% 纲要位置：几何学计划纲要.md，第 1440 行。
% 正文状态：初稿；尚未正式审核。

% 纲要内容开始
% 
% * 杯积和 Leibniz 规则
% * 边缘映射的几何意义
% * 传递映射的代表元解释
% * 映射纤维化之间的比较
% 
% 纲要内容结束

\section{乘积、传递与自然性}
\label{g2:sec:1102}

第二页给出了许多群，真正的计算还需要它们之间的联系。乘法使一个生成元上的微分决定整片表格；边缘映射把表格接回纤维包含与基底投影；传递则把纤维上的一个类转为基底上的障碍。下面始终使用含幺交换系数环 \(R\)，谈到零行的普通系数时另假设纤维道路连通。

\subsection{杯积和 Leibniz 规则}
\label{g2:subsec:1102:01}

纤维的杯积与沿路运输交换，所以 \(\mathcal H^*(F;R)\) 是分次代数的局部系统。在局部平凡的记法中，\(a\otimes u\) 表示一个基底 \(p\) 次类配上一个纤维 \(q\) 次类。将两个这样的类相乘，必须先把第一个纤维因子越过第二个基底因子。因此第二页的乘法为
\begin{equation}\label{g2:eq:serre-product-sign}
	(a\otimes u)(b\otimes v)=(-1)^{q s}(a\smile b)\otimes(u\smile v),\qquad |a|=p,\ |u|=q,\ |b|=s.
\end{equation}
这是有张量识别时的公式。一般局部系数情形，则先按第5章的道路运输把系数移到同一纤维，再作纤维杯积；同一个 \((-1)^{qs}\) 仍须保留。

\begin{theorem}[Serre 谱序列的乘法]\label{g2:thm:serre-multiplicative}
	上同调 Serre 谱序列从第二页起是按总次数分次交换的代数谱序列。对 \(x\in E_r^{p,q}\)、\(y\in E_r^{s,t}\)，有
	\begin{equation}\label{g2:eq:serre-leibniz}
		\mathrm d_r(xy)=(\mathrm d_rx)y+(-1)^{p+q}x\,\mathrm d_ry.
	\end{equation}
	下一页的乘法由本页乘法诱导；第二页乘法由局部系数杯积和\cref{g2:eq:serre-product-sign}给出。目标上同调的滤过满足 \(F^pH^m(E)F^sH^n(E)\subset F^{p+s}H^{m+n}(E)\)，且极限页作为代数同构于这个环的伴随分次。
\end{theorem}
\begin{proof}
	先把几何滤过落实为一个可作乘法的链模型。用\cref{g2:lem:serre-filtered-cw-model}的相容 CW 逼近 \(X_p\to E_p\)，记 \(X=\bigcup_pX_p\)、\(C_*=C_*^{\mathrm{cell}}(X;\mathbb Z)\)。给每个胞腔记录它首次出现的层数；\(C_*(X_p)\) 即由层数不超过 \(p\) 的胞腔生成。积空间使用胞腔积的弱拓扑，积胞腔的层数是两因子层数之和。第4章的交叉积将其胞腔链识别为 \(C_*\otimes C_*\)，边界为 \(\partial(c\otimes d)=\partial c\otimes d+(-1)^{|c|}c\otimes\partial d\)。

	取基底对角映射的胞腔逼近 \(\Delta_c:B\to B\times B\)。它满足
	\[
		\Delta_c(B^k)\subset(B\times B)^k=\bigcup_{i+j\le k}B^i\times B^j.
	\]
	在 Hurewicz 替换中，将 \(\Delta_B\simeq\Delta_c\) 的同伦从总空间的真实对角映射开始提升，得到 \(\widetilde\Delta:E\to E\times E\)，且 \(\widetilde\Delta(E_k)\subset\bigcup_{i+j\le k}E_i\times E_j\)。在上述 CW 模型上逐层比较并作胞腔逼近，得到链映射
	\[
		D:C_*\longrightarrow C_*\otimes C_*,\qquad
		D(C_*(X_k))\subset\sum_{i+j\le k}C_*(X_i)\otimes C_*(X_j).
	\]
	这里逐层比较的依据是：相容 CW 逼近在每层都是弱同伦等价；对一个待延拓的胞腔，比较映射及比较同伦的障碍是该弱等价的一个相对同伦类，因而为零。按胞腔维数延拓，并在已处理子复形上固定选择，就得到相容比较。对有限个积因子的并也同样逐个积胞腔处理。这说明所用链对角确实保留总层数，而不只是普通上同调上的某个对角。

	令 \(C^*=\operatorname{Hom}(C_*,R)\)，滤过取在 \(C_*(X_{p-1})\) 上为零的上链。定义 \(uv=(u\otimes v)D\)。若 \(u\in F^p\)、\(v\in F^s\)，对层数小于 \(p+s\) 的任一积胞腔，至少有一个因子被相应上链消去，所以 \(uv\in F^{p+s}\)。\(D\) 为链映射又给出严格的上链恒等式
	\[
		\mathrm d(uv)=(\mathrm du)v+(-1)^{|u|}u\,\mathrm dv.
	\]
	若 \(u,v\) 是第 \(r\) 页的代表，即 \(\mathrm du\in F^{p+r}\)、\(\mathrm dv\in F^{s+r}\)，右端属于 \(F^{p+s+r}\)。更换代表时，乘法与上式使新增项落入第10章商公式的分母：多一层的近似闭元仍多一层，上边界乘积则改写为一个上边界加一个更高滤过项。因此乘积良定义，并且直接取上式在目标商中的类便得到\cref{g2:eq:serre-leibniz}。此论证对每个 \(r\) 适用，无需从第一微分的导子性质猜测后续微分。

	还要说明结合律与交换律为何在第二页成立。对角的两种三重迭代、交换两因子的对角，以及不同胞腔逼近，均由真实对角提供同伦。把同伦作相对胞腔逼近时，\(B^k\times I\) 只需进入总骨架 \(k+1\)。故所得链同伦提高层数至多一；对偶的上链同伦 \(K\) 降低滤过至多一。对于第 \(r\ge2\) 页的闭代表，总滤过为 \(P\) 时，误差形如 \(\mathrm dK(w)+K(\mathrm dw)\)。第一项是来自 \(F^{P-1}\subset F^{P-r+1}\) 的允许边界，第二项属于 \(F^{P+r-1}\subset F^{P+1}\)，两者在该页均为零。这证明从第二页起的结合律、单位律和总次数交换律。链对角代表真实对角，所以在总上同调上得到通常杯积。

	最后，在一个 \(p\) 胞腔与一个 \(s\) 胞腔上，用基底相对定向类乘纤维类识别第一页。交换中间的 \(q\) 次纤维因子和 \(s\) 次基底因子给出 \((-1)^{qs}\)，故第二页为所述乘法。对真闭代表同样取乘积，所得滤过和极限页乘积便是目标杯积的伴随分次。
\end{proof}

计算时需要区分两个次数：乘积交叉符号是 \((-1)^{qs}\)，导子公式的符号是 \((-1)^{p+q}\)。后者使用整个类的总次数。

\subsection{边缘映射的几何意义}
\label{g2:subsec:1102:02}

设 \(F\xrightarrow{i}E\xrightarrow{p}B\)，其中 \(B,F\) 道路连通。第二页零行是 \(H^*(B;R)\)，零列是纤维上同调中沿全部闭路运输不变的部分：
\[
	E_2^{n,0}=H^n(B;R),\qquad E_2^{0,n}=H^n(F;R)^{\pi_1(B)}.
\]
零行上的类只能被别的类击中，零列上的类只能向外发出微分。因此前者逐页取商，后者逐页取子群。

\begin{proposition}[边缘映射]\label{g2:prop:serre-edge-maps}
	Serre 谱序列中的两条边缘分别分解投影拉回与纤维限制：
	\[
		H^n(B;R)\twoheadrightarrow E_\infty^{n,0}=F^nH^n(E;R)\hookrightarrow H^n(E;R),
	\]
	\[
		H^n(E;R)\twoheadrightarrow E_\infty^{0,n}\hookrightarrow H^n(F;R)^{\pi_1(B)}\hookrightarrow H^n(F;R).
	\]
	两条复合依次为 \(p^*\) 和 \(i^*\)。特别地，能从总空间限制下来的纤维类必为永久圈。
\end{proposition}
\begin{proof}
	把原纤维化映到恒等纤维化 \(B\to B\)。后者仅有零行，谱序列在第二页退化，诱导的总上同调映射为 \(p^*\)；对骨架对的拉回以及相应正合偶取映射，即得第一条分解。由于总次数 \(n\) 中 \(F^{n+1}=0\)，最右的分级商就是实际子群 \(F^n\)。

	再从基点上的纤维化 \(F\to\{b_0\}\) 映入原纤维化。零骨架上的限制把全局类送到各纤维上的相容类，沿边的相容性正是运输不变性。滤过的定义给出 \(F^1H^n(E)=\ker i^*\)：若基底有多个顶点，沿路径运输表明在一个纤维为零等价于在全部顶点纤维为零。于是 \(H^n(E)/F^1\) 就是限制像，得到第二条分解。
\end{proof}

\begin{corollary}\label{g2:cor:serre-section-bottom-row}
	若 \(p\) 有截面，则所有落入零行的微分均为零。
\end{corollary}
\begin{proof}
	截面 \(s\) 满足 \(s^*p^*=\operatorname{id}\)，所以 \(p^*\) 单射。由第一条边缘分解，零行任何类都不能在逐页取商时被消去。
\end{proof}

这个判据常比直接追踪代表元快捷。它控制的是零行；其他行之间仍可能存在微分。

\subsection{传递映射的代表元解释}
\label{g2:subsec:1102:03}

设 \(q>0\)。一个纤维类先要保持运输不变，才能进入 \(E_2^{0,q}\)；随后还要通过 \(\mathrm d_2,\ldots,\mathrm d_q\) 的检验，才能到达第 \(q+1\) 页。最后可能发出的箭头恰为
\begin{equation}\label{g2:eq:serre-transgression}
	\tau=\mathrm d_{q+1}:E_{q+1}^{0,q}\longrightarrow E_{q+1}^{q+1,0}.
\end{equation}
它称为\term[chuandi]{传递}[transgression]。源是 \(H^q(F;R)\) 的一个子群，目标是 \(H^{q+1}(B;R)\) 的一个商；较早微分可能已在零行中加入关系。符号 \(\tau a\in H^{q+1}(B;R)\) 只有在该商确实等于整个基底群时才直接使用。

\begin{proposition}[传递的延拓解释]\label{g2:prop:serre-transgression}
	类 \(a\in H^q(F;R)^{\pi_1(B)}\) 属于 \(E_{q+1}^{0,q}\)，当且仅当它可延拓为 \(E_q=p^{-1}(B^q)\) 上的 \(q\) 次上同调类。选择这样的延拓后，它在 \((q+1)\) 胞腔上继续延拓的障碍组成一个基底上链；该上链所定义的类，模去较早微分的像，就是 \(\tau a\)。
\end{proposition}
\begin{proof}
	先对限制纤维化 \(E_q\to B^q\) 使用上一命题。零列 \((0,q)\) 的所有可能非零出射在第 \(q\) 页以前已经出现；其余箭头落在基底维数以外。因此这个位置的永久圈恰为 \(E_{q+1}^{0,q}\)，而边缘命题把它识别为 \(H^q(E_q;R)\to H^q(F;R)\) 的像。

	用一个闭上链 \(A\) 表示所选延拓，再将其作为上链延拓到 \(E_{q+1}\)。余边界 \(\mathrm d\widetilde A\) 在 \(E_q\) 上为零，故定义
	\[
		[\mathrm d\widetilde A]\in H^{q+1}(E_{q+1},E_q;R)
		\cong C^{q+1}(B;\mathcal H^0(F;R))=C^{q+1}(B;R).
	\]
	这个等式描述每个新胞腔上的障碍值。再取一次余边界为零，故这些值组成胞腔上圈。更换上链延拓只改动相对边界；更换较低骨架上的修正，则加入前面几页已经允许的边界。第10章代表元公式把这些不确定性恰好除去，并把剩余的 \(\mathrm d\widetilde A\) 识别为 \(\mathrm d_{q+1}[a]\)。
\end{proof}

还可用空间对表达一种十分实用的判定。记 \(\delta:H^q(F;R)\to H^{q+1}(E,F;R)\) 为连接同态。若 \(b\in H^{q+1}(B,b_0;R)\) 满足
\begin{equation}\label{g2:eq:transgression-relative-criterion}
	\delta a=p^*b,
\end{equation}
则 \(a\) 可传递，且 \(\tau a\) 是 \(b\) 在 \(E_{q+1}^{q+1,0}\) 中的像。事实上，在相容胞腔模型中取代表并调整一个相对上边界，就能使 \(\mathrm dA=p^*b\)；右端在基底 \(q\) 骨架上为零，于是 \(A\) 先给出上述延拓，再在下一层留下 \(b\)。连接同态采用“延拓后取余边界”的约定，因此这里没有额外负号。

\ProseParagraphBreak

例如，若 \(\widetilde H^j(F;R)=0\) 对 \(0<j<q\) 成立，则从 \((0,q)\) 出发的较早微分及落入 \((q+1,0)\) 的较早微分都没有非零源或靶。此时所有运输不变的 \(q\) 次类都可传递，目标也是整个 \(H^{q+1}(B;R)\)。球丛的特征类正是由这种情形产生。

\subsection{映射纤维化之间的比较}
\label{g2:subsec:1102:04}

给定交换图
\[
	\begin{tikzcd}
		F\arrow[r]\arrow[d,"f_F"]&E\arrow[r,"p"]\arrow[d,"f_E"]&B\arrow[d,"f_B"]\\
		F'\arrow[r]&E'\arrow[r,"p'"]&B',
	\end{tikzcd}
\]
沿路提升与 \(f_E\) 相容到同伦，因此纤维上同调的拉回给出局部系统映射 \(f_B^*\mathcal H^q(F';R)\to\mathcal H^q(F;R)\)。这才是第二页映射的完整输入；当基本群作用不平凡时，只列两个抽象群之间的同态还不够。

\begin{proposition}[自然性与比较]\label{g2:prop:serre-comparison}
	上述图诱导反向的谱序列映射 \(f_r^*:E_r'{}^{p,q}\to E_r^{p,q}\)，与微分、乘法及传递交换；第二页为局部系数上同调上的诱导映射，极限页为滤过映射 \(f_E^*\) 的伴随分次。如果第二页在所有位置都是同构，则 \(f_E^*:H^*(E';R)\to H^*(E;R)\) 是同构。
\end{proposition}
\begin{proof}
	先将 \(f_B\) 胞腔逼近，并在 Hurewicz 替换中提升这段同伦以相应调整 \(f_E\)。所得映射保持骨架逆像，故给出空间对正合列之间的交换图。第10章的导出正合偶构造于是逐页给出 \(f_r^*\)，且与微分交换。逐胞腔的相对群识别使用纤维运输，所以第一微分及第二页映射恰好是所述局部系数映射。对角映射的自然性使乘积相容；传递本身是一条指定微分，故也相容。

	若 \(f_2^*\) 为同构，取各页上同调归纳得到所有 \(f_r^*\) 为同构。固定总次数 \(n\)，两边的目标滤过只有 \(n+1\) 层。从最深一层开始，将相邻滤过项及其商放入短正合列，逐层用五引理，便得到 \(f_E^*\) 在 \(n\) 次的同构。这里使用的是前节已经证明的有限目标滤过。
\end{proof}

对同调版也有同样的比较，箭头方向与连续映射一致。特别地，若基底均单连通，\(f_B\) 与 \(f_F\) 都诱导整数同调同构，则 \(f_E\) 也诱导整数同调同构：第4章的同调万有系数短正合列先将基底同构提升为任意纤维同调系数下的同构，随后逐页比较即可。这个步骤解释了为什么“基底与纤维分别不变”能够推出总空间的同调不变。
% 纲要条目排为编号小节标题，并进入目录与书签；说明段仍保留于纲要注释中。
